\begin{table}[!htb]\centering\footnotesize\setlength{\tabcolsep}{3.2pt}
\caption{Per-revolution statistics computed from the manuals.}\label{tab:stats}
\begin{tabular}{@{}lrrrrrrrrrrrrr@{}}\toprule
Rev. & terms & tok. & E & P & T & N & inf. & H & R & X & U & $\tau$ & $\kappa$\\\midrule
CHEM & 23 & 35 & 8 & 11 & 4 & 0 & 31 & 20 & 0 & 2 & 9 & 0.879 & 0.710\\
ASTR & 25 & 35 & 6 & 10 & 6 & 3 & 34 & 16 & 6 & 6 & 6 & 0.887 & 0.824\\
MECH & 23 & 37 & 3 & 13 & 6 & 1 & 35 & 12 & 13 & 1 & 9 & 0.845 & 0.743\\
HEAT & 23 & 30 & 8 & 12 & 1 & 2 & 30 & 16 & 3 & 2 & 9 & 0.897 & 0.700\\
QUAN & 21 & 27 & 2 & 16 & 3 & 0 & 30 & 11 & 7 & 7 & 5 & 0.925 & 0.833\\
BIOL & 24 & 28 & 10 & 10 & 3 & 1 & 31 & 18 & 0 & 8 & 5 & 0.908 & 0.839\\
PHYS & 26 & 31 & 3 & 15 & 4 & 4 & 30 & 13 & 0 & 6 & 11 & 0.783 & 0.633\\
OPTI & 18 & 21 & 8 & 7 & 2 & 1 & 25 & 15 & 1 & 2 & 7 & 0.812 & 0.720\\
MIAS & 21 & 29 & 6 & 8 & 5 & 2 & 31 & 17 & 1 & 2 & 11 & 0.793 & 0.645\\
HERE & 17 & 21 & 6 & 7 & 2 & 2 & 23 & 11 & 3 & 5 & 4 & 0.895 & 0.826\\
\midrule
All & 221 & 294 & 60 & 109 & 36 & 16 & 300 & 149 & 34 & 41 & 76 & & \\\bottomrule\end{tabular}\end{table}

\begin{table}[!htb]\centering\footnotesize
\caption{Further indices: unweighted translation share, inferential preservation $\pi$, disagreement share $\delta$, the untranslatable cluster and the neutral core.}\label{tab:further}
\begin{tabular}{@{}lrrrp{6.2cm}r@{}}\toprule
Rev. & $\tau_{\mathrm{unw}}$ & $\pi$ & $\delta$ & untranslatable cluster $C$ & core\\\midrule
CHEM & 0.931 & 0.645 & 0.065 & CH01 (phlogiston), CH02 (dephlogisticated air), CH03 (phlogisticated air) & 3\\
ASTR & 0.840 & 0.647 & 0.176 & AS03 (primum mobile), AS04 (planetary orb), AS08 (equant), AS15 (centre of the universe), AS24 (aether) & 3\\
MECH & 0.819 & 0.714 & 0.029 & ME02 (force (impressed)), ME06 (absolute space), ME07 (absolute time), ME16 (ether), ME19 (weight), ME22 (rigid body), ME23 (action at a distance) & 2\\
HEAT & 0.891 & 0.633 & 0.067 & HE01 (caloric), HE06 (free caloric), HE07 (combined caloric) & 1\\
QUAN & 0.937 & 0.600 & 0.233 & QU01 (trajectory), QU02 (electron orbit), QU18 (state (mechanical)) & 0\\
BIOL & 0.896 & 0.581 & 0.258 & BI05 (plan of creation), BI06 (archetype), BI10 (centre of creation), BI16 (chain of being) & 7\\
PHYS & 0.769 & 0.433 & 0.200 & GA03 (vital spirit), GA04 (natural spirit), GA05 (animal spirit), GA08 (septal pores), GA14 (innate heat), GA15 (sooty vapours), GA18 (faculties), GA26 (rete mirabile) & 2\\
OPTI & 0.889 & 0.640 & 0.080 & OP02 (light corpuscle), OP03 (fits of easy reflection and transmission), OP08 (sides of rays) & 4\\
MIAS & 0.810 & 0.581 & 0.065 & MI01 (miasma), MI03 (ferment), MI05 (contagium), MI07 (spontaneous generation), MI08 (epidemic constitution), MI12 (sewer gas), MI13 (cholera poison) & 5\\
HERE & 0.824 & 0.609 & 0.217 & HR01 (gemmule), HR05 (use-inheritance), HR06 (telegony), HR13 (ancestral heredity) & 3\\
\bottomrule\end{tabular}\end{table}

\begin{table}[!htb]\centering\footnotesize
\caption{Translation index under the standard scoring S0 and the alternative schemes S1--S4.}\label{tab:sens}
\begin{tabular}{@{}lrrrrrc@{}}\toprule
Rev. & S0 & S1 & S2 & S3 & S4 & C8\\\midrule
CHEM & 0.879 & 0.741 & 1.000 & 0.588 & 0.931 & pass\\
ASTR & 0.887 & 0.560 & 0.920 & 0.667 & 0.840 & pass\\
MECH & 0.845 & 0.724 & 0.983 & 0.509 & 0.819 & pass\\
HEAT & 0.897 & 0.838 & 0.956 & 0.566 & 0.891 & pass\\
QUAN & 0.925 & 0.828 & 1.000 & 0.529 & 0.937 & pass\\
BIOL & 0.908 & 0.837 & 0.980 & 0.633 & 0.896 & pass\\
PHYS & 0.783 & 0.717 & 0.850 & 0.467 & 0.769 & pass\\
OPTI & 0.812 & 0.750 & 0.875 & 0.604 & 0.889 & pass\\
MIAS & 0.793 & 0.556 & 0.933 & 0.615 & 0.810 & pass\\
HERE & 0.895 & 0.842 & 0.947 & 0.566 & 0.824 & pass\\
\bottomrule\end{tabular}\end{table}

\begin{table}[!htb]\centering\footnotesize
\caption{Outcome of each revolution on the five shared criteria ($+$ favoured the successor, $-$ the predecessor, $\pm$ mixed or contested, ? no claim made), with the mechanical checks C4 and C5.}\label{tab:crit}
\begin{tabular}{@{}lcccccccc@{}}\toprule
Rev. & acc. & cons. & scope & simp. & fruit. & C4: $0<\tau<1$ & C5: some $+$ & no $-$\\\midrule
CHEM & $+$ & ? & $\pm$ & ? & $+$ & pass (0.879) & pass (2) & yes\\
ASTR & $+$ & $-$ & $+$ & $+$ & $+$ & pass (0.887) & pass (4) & no\\
MECH & $+$ & $+$ & $+$ & $\pm$ & $+$ & pass (0.845) & pass (4) & yes\\
HEAT & $+$ & ? & $+$ & ? & $+$ & pass (0.897) & pass (3) & yes\\
QUAN & $+$ & $+$ & $+$ & ? & $+$ & pass (0.925) & pass (4) & yes\\
BIOL & ? & $-$ & $+$ & $+$ & $+$ & pass (0.908) & pass (3) & no\\
PHYS & $+$ & $+$ & $\pm$ & ? & $+$ & pass (0.783) & pass (3) & yes\\
OPTI & $+$ & ? & $+$ & ? & $+$ & pass (0.812) & pass (3) & yes\\
MIAS & $+$ & ? & $+$ & ? & $+$ & pass (0.793) & pass (3) & yes\\
HERE & $+$ & ? & $+$ & ? & $+$ & pass (0.895) & pass (3) & yes\\
\bottomrule\end{tabular}\end{table}

\begin{crittable}
CHEM & accuracy & $+$ & Gravimetric balance: the weight gained in calcination equals the weight of air absorbed (Lavoisier, 1770s); the phlogistic accounts of the gain were not quantitative.\\
CHEM & scope & $\pm$ & The oxygen theory explained combustion, calcination and respiration together; the phlogistic theory explained why metals are alike and their calxes unlike (Kuhn; Chang).\\
CHEM & fruitfulness & $+$ & The new nomenclature (1787), the conservation of weight as a working principle, and quantitative analysis, on which Dalton built.\\
\midrule
ASTR & accuracy & $+$ & Kepler's Rudolphine Tables (1627) were much more accurate than earlier tables; Kepler predicted the transit of Mercury of 1631, observed by Gassendi. (Copernicus's own tables were not more accurate than Ptolemaic ones; Kuhn 1957, 1977.)\\
ASTR & consistency & $-$ & External consistency with Aristotelian physics and cosmology favoured the geocentric system until the new mechanics (Kuhn 1977).\\
ASTR & scope & $+$ & Fixed the order and relative distances of the planets from observation and explained the link of the second anomaly to the Sun.\\
ASTR & simplicity & $+$ & With Kepler, one ellipse per planet governed by the area law replaced compounded circles and equants.\\
ASTR & fruitfulness & $+$ & Led to Kepler's laws and, through them, to Newton's mechanics.\\
\midrule
MECH & accuracy & $+$ & General relativity accounted for the anomalous 43 arcseconds per century of Mercury's perihelion advance and predicted the deflection of starlight reported in 1919.\\
MECH & consistency & $+$ & Special relativity removed the conflict between the Galilean invariance of mechanics and the Lorentz invariance of Maxwell's equations.\\
MECH & scope & $+$ & Mechanics, electrodynamics and optics of moving bodies brought under one kinematics; gravitation and inertia unified.\\
MECH & simplicity & $\pm$ & Conceptual unity was gained; the mathematics of general relativity is far more complex.\\
MECH & fruitfulness & $+$ & Mass--energy equivalence, relativistic cosmology, and the relativistic theories of the electron and of fields.\\
\midrule
HEAT & accuracy & $+$ & Joule's determinations of the mechanical equivalent of heat (1840s) showed a fixed ratio of work spent to heat produced, which the conservation of caloric excluded.\\
HEAT & scope & $+$ & The kinetic theory derived the gas laws and predicted that the viscosity of a gas is independent of its pressure (Maxwell, 1860), which Maxwell confirmed experimentally (1866).\\
HEAT & fruitfulness & $+$ & The first and second laws, statistical mechanics, and the theory of transport processes.\\
\midrule
QUAN & accuracy & $+$ & Planck's law (1900) fitted the black-body spectrum where the Rayleigh--Jeans law failed; Einstein's theory of specific heats (1907) explained their fall at low temperature; the hydrogen spectrum was derived.\\
QUAN & consistency & $+$ & The classical nuclear atom is unstable by classical electrodynamics; quantum mechanics yields stable stationary states.\\
QUAN & scope & $+$ & Chemical bonding, the periodic table, spectra and the solid state.\\
QUAN & fruitfulness & $+$ & Quantum chemistry, solid-state physics, nuclear and particle physics.\\
\midrule
BIOL & consistency & $-$ & External: Kelvin's estimates of the age of the Earth (1860s) seemed to leave too little time, and Jenkin (1867) argued that blending inheritance would swamp favourable variations.\\
BIOL & scope & $+$ & One hypothesis explained the natural system, homology, rudimentary organs, embryonic resemblance, the succession of types and geographical distribution (Origin, chs.~11--14).\\
BIOL & simplicity & $+$ & One process of descent with modification replaced a separate act of creation for each species.\\
BIOL & fruitfulness & $+$ & Phylogenetic systematics, biogeography and, after the synthesis with genetics, population biology.\\
\midrule
PHYS & accuracy & $+$ & Harvey's quantitative argument: the blood expelled by the heart in a short time far exceeds what food could supply, so the same blood must return (De motu cordis, ch.~9).\\
PHYS & consistency & $+$ & The circulation required no invisible septal pores, whose existence Vesalius had come to doubt in the 1555 edition of the Fabrica.\\
PHYS & scope & $\pm$ & Harvey explained the motion of the blood and the function of the venous valves; the Galenic system offered accounts of nutrition, respiration and heat that the circulation left open.\\
PHYS & fruitfulness & $+$ & Malpighi's capillaries (1661), the Oxford physiologists' work on respiration, transfusion experiments.\\
\midrule
OPTI & accuracy & $+$ & Fresnel's integrals gave the positions and intensities of diffraction fringes; Arago observed the bright spot at the centre of a disk's shadow that Poisson had derived from Fresnel's theory (1819).\\
OPTI & scope & $+$ & Interference, diffraction, polarization and double refraction brought under transverse waves.\\
OPTI & fruitfulness & $+$ & Fresnel's equations for reflection and refraction; the electromagnetic theory of light.\\
\midrule
MIAS & accuracy & $+$ & Specific organisms were isolated and shown to cause specific diseases (Koch on anthrax, 1876; tuberculosis, 1882; the cholera vibrio, 1884), and the routes of transmission were traced.\\
MIAS & scope & $+$ & Fermentation, putrefaction, wound infection and the specificity of epidemic diseases explained by one hypothesis.\\
MIAS & fruitfulness & $+$ & Antiseptic surgery (Lister, 1867), attenuated vaccines (Pasteur on anthrax, 1881; rabies, 1885).\\
\midrule
HERE & accuracy & $+$ & The Mendelian ratios (3:1, 9:3:3:1) were confirmed in crosses after 1900 (Correns, de Vries, Tschermak, Bateson).\\
HERE & scope & $+$ & Segregation, independent assortment, linkage and sex-linked inheritance (Morgan, 1910) brought under one chromosomal mechanism.\\
HERE & fruitfulness & $+$ & The Hardy--Weinberg principle (1908) and population genetics; Fisher (1918) derived the biometric correlations from Mendelian factors.\\
\bottomrule
\end{crittable}
